\documentclass[year=2026]{sit}

\usepackage{graphicx}
\usepackage{booktabs}
\usepackage[normalem]{ulem}
\usepackage{siunitx}
\sisetup{%
  group-digits=integer,%
  group-separator={,}%
}

\paperid{99999}

\title{講演論文作成について}
\etitle{Making Research Paper}

\author{%
  機械 太郎\af{1}，\presenter{技術 さくら\af{2}}\\%
  機械 二郎\af{1}，機械 三郎\af{1}，東京 花子\af{3}%
}
\eauthor{%
  Taroh KIKAI\af{1}, Sakura GIJYUTSU\af{2},\\%
  Jiroh KIKAI\af{1}, Saburoh KIKAI\af{1} and Hanako TOKYO\af{3}
}

\affiliation{%
  \afdef{1}{日本機械大学}{Nihon Kikai University}
  \afdef{2}{信濃町大学}{Shinanomachi University}
  \afdef{3}{機械株式会社}{Kikai Corporation}
}

\abstract{%
When preparing the manuscript, read and observe this sample carefully. This sample was prepared using \LaTeX. Please provide an abstract here in about 150--250 words. ...............................................................................................
.......................................................................................................................................................................................
.......................................................................................................................................................................................
.......................................................................................................................................................................................
.......................................................................................................................................................................................
.............................................................................................
}

\keywords{Keyword 1, Keyword 2, Keyword 3, ... (Provide five to ten keywords.)}

\begin{document}

\maketitle


\section{諸　　　言}\label{sec:introduction}

このテンプレートファイルは，Sports Informatics and Technology 2025（SIT2025）の講演論文原稿を準備するにあたり，原稿体裁を整えて作成することができるようにスタイルファイルとして，MS-Word形式でのテンプレートを参考に作成したものである．
配布されているMS-Word形式のテンプレートに準拠したフォーマットになっている．

本文の文字数はMS-Word形式のテンプレートに準拠し，1ページ当たり，50文字×46行×1段組で2300字としている．
なお，配布されているMS-Word形式のテンプレートファイルを用いた場合と，本テンプレートを使用した場合とで仕上がりに差異が生じるが，設定上は条件を守っている．
文章の区切りには全角の読点「，」（カンマ）と句点「．」（ピリオド）を用いる．カッコも全角入力する．

コンパイラは LuaLaTeX を対象としており，他のコンパイラ（pdfLaTeX や upLaTeX，XeLaTeX）ではエラーにより正常に出力されない場合がある．

SIT2025の講演論文原稿では，PDF化した提出原稿のファイルサイズを，１ファイルあたり5 MB以下とする．ファイル名は，講演発表申込時の確認メールに記載された『申込番号』の下3桁を使用する．例えば，申込番号がU00555の場合は，555.pdf とすること．なお，申込番号が1桁あるいは2桁の場合は，001pdfや067.pdfのように0を補うこと．


\section{このテンプレートファイルの使い方}\label{sec:usage}

このテンプレートの表題（副題），著者名，本文などはあらかじめ指定のフォントサイズなどの書式が設定されている．
特に設定をいじらなければ，文字数，行数など定められた体裁で論文を作成することができる．
しかし，絶対的な出来上がりのレベルを保証するものではないので，体裁が望むレベルに達しない場合には，使用の環境に合わせ，筆者各自において微調整を行うなど，論文集掲載の体裁に最も近い設定を行う必要がある．

なお，書式の変更には .cls ファイルおよび jlreq クラスへの理解が必要になる．


\section{原稿執筆の手引き}\label{sec:draft}

\subsection{原稿の規定ページ数について}
SIT 2026における原稿のページ数は，2ページを標準とするが，最大で4ページまで可とする．
原稿中にページ番号は不要である． 

\subsection{原稿の作成に際して}
原稿の冒頭には，和文の表題・副題，著者名，英文の表題・副題，ローマ字著者名（苗字は例に従い大文字），ローマ字連絡先所属機関（英文の著者連絡先には，代表者の所属機関名・部署名等）を入れる．

\subsection{表題及び副題の付け方}
原稿の表題は内容を明確に表現するもので，しかも簡潔なものが望まれる．
また，必要に応じて副題を付けてもよいらしいが，このテンプレートファイルは現在のところ副題に対応していない．
なお，副題をつける場合には両括弧でくくる．

\subsection{英文抄録の書き方}
長さは150～250語程度で，途中で改行をしないで，本文と切り離してそれだけを読んでも，論文の内容が具体的に分かるように研究対象，研究方法・装置，結果について書く．また，本文中の図・表・文献は，引用しない．式を書く必要がある場合は，式の番号を引用せずに，式をそのまま書く．

\subsection{キーワードの付け方}
キーワードは，論文の内容を代表する重要な用語である．
これによって論文の分類，検索が迅速になる．選定に際しては，基準キーワード集の選定要領に従い選ぶ．
\begin{itemize}
  \item[(1)] キーワードは，5～10語句とする．
  \item[(2)] キーワードは，英文抄録の直下に英語で記載する．キーワードにはハイフンは使用せず（ハイフンを使用してひとつの単語として一般に認識されるものを除く），前置詞や冠詞も含めない．
\end{itemize}

\subsection{見出し（章，節，項）の付け方及び書き方}
本文は適当に区分して，見出しを付ける．
体裁は \LaTeX が良い感じにしてくれるので，余計な改行を入れるなどの処理は気にしなくてよい．
なお，現在のところ \verb|\section| および \verb|\subsection| のみに対応している．

\subsection{量記号・単位記号の書き方}
量記号はイタリック体，単位記号はローマン体とする．無次元数はイタリック体で書く．

\subsection{用いる単位について}
単位は，SI単位を使用する．数学記号・単位記号及び量記号は，半角英数字を使用する．
なお，SI単位については，日本機械学会発行の「機械工学SIマニュアル」および「JISZ8203 国際単位系（SI）及びその使い方」を参照する．
記載時には，siunitx パッケージを使用すると良い．

\subsection{用いる記号}
数学記号は，JISZ8201に従う．また，量を表す文字記号（量記号）は，JISZ8202に従う．
詳細は MS-Word 形式のテンプレートファイルを参照のこと．


\section{図及び写真・表の作成に関して}\label{sec:figures_and_tables}
\begin{itemize}
  \item[（1）] 本文中では，図1，表1のように日本語で書く．写真は，図として扱う．カラーで掲載できる．
  \item[（2）] 番号・説明などは，図についてはその下に，表についてはその上に書く．
  \item[（3）] \sout{本文と，図・表の間は1行以上の空白を空けて，見やすくする．}
  \item[（4）] 図中・表中の説明及び題目はすべて英語で書く（最初の文字は大文字とする）．例に従い該当図表が示す重要知見に係る説明もあわせて記述すること．図のCaptionの書体としてはSerif系を利用し，9.5ポイントの大きさで記載する．
  \item[（5）] 図及び表の横に空白ができても，その空白部には本文を記入してはならない．
  \item[（6）] 図及び表は，余白部分にはみ出してはならない．
  \item[（7）] 論文の内容をわかりやすく的確に伝えるために，図・表を適当に用いて，本文と対応するページにレイアウトする．
\end{itemize}

\begin{table}[tb]
  \begin{center}
    \caption{Sample table}\label{tab:sample}
    \begin{tabular}{cc}
      \toprule
      Category & size\\
      \midrule
      Train & \num{100000}\\
      \bottomrule
    \end{tabular}
  \end{center}
\end{table}


\section{数式の書き方}\label{sec:equations}
\sout{式番号は，式と同じ行に右寄せして（ ）の中に書く．
また，}本文で式を引用するときは，式（1）のように書く．
\sout{式を書くときは，2文字分空白を空ける．}
また，必要行数分を必ず使うようにして書く．
3行必要とする式を2行につめて書いたり，2行に分かれる式を1行に収めたりしない．
なお，本文と式，式相互間は1行以上の空白を空けて，見やすくする．
また，原則として数式エディタのポイント数は本文に準じるものとするが，添え字等が小さく読みにくくなるときは適宜拡大する．

\begin{align}
  & d\left\{\sum \frac{1}{2}mk\left[%
    \left(\frac{dx_i}{dt}\right)^2%
    + \left(\frac{dy_i}{dt}\right)^2%
    + \left(\frac{dz_i}{dt}\right)^2%
  \right]\right\} = \sum\left(X_i dx_i + Y_i dy_i + Z_i dz_i\right)\\
  & \overline{C}(t) = \frac{1}{N}\sum_{i=1}^{N} C_i(t)\\
  & C_\mathit{th} = \frac{\displaystyle\sum_{i=1}^N C(t)}{\sqrt{\displaystyle\sum_{i=0} (\ell_0 - \ell_1)}} = \frac{b}{a}C_0
\end{align}

数式の上下にはこのくらいのスペースがあく．


\section{引用文献の書き方}\label{sec:citation}

他者の報告・データなどを引用するときには，必ずその出所を明示しなければならない．
また，研究の背景を説明する際には，必ず必要な文献を引用する．引用文献の書き方の詳細については，以下の執筆要綱およびテンプレートの記載例を参照すること．

\begin{itemize}
  \item[（1）] 一般に公表されていない文献（投稿予定および投稿中で未発表の論文も含む）は，できるだけ引用しない．また，Webサイトの引用については止むを得ない場合を除き，できるだけ引用しない．引用せざるを得ない場合は，必ず参照日付を明記する．（例  文献 (参照日 2013年8月14日)）
  \item[（2）] 本文中の引用箇所には，著者名と発行年を記載する．（日本語文献例　(竹内，2005)　英語文献例 (Ahrendt and Taplin, 1951)）3名以上の著者がいる場合の著者名の記載方法は，代表著者名他の記載とする．（日本語文献例 (蔦原他，2003)  英語文献例 (Takeuchi, et al., 2006)）発行年が同じである同じ著者からの2つ以上の引用を記載する場合には，発行年の後にa,b,cを記載する．（例　(Karin and Hanamura, 2010a, 2010b)）
  \item[（3）] 引用した文献は，本文末尾に文献としてアルファベット順にまとめて記載する．
  \item[（4）] 文献に記載する誌名は，略記せずにすべてを記載する．
  \item[（5）] 日本語の文献を引用する場合は日本語表記とし，英語の文献を引用する場合は英語表記とする．
  \item[（6）] bibtex を使用することで \cite{shinjuku:1987,keer:1986} のように出力される．
\end{itemize}


\section{結　　　言}\label{sec:conclusion}
学会などの原稿テンプレートは，MS-Wordファイルなどという研究者を不幸にするだけの形式ではなく，最初から \LaTeX 形式で配布して欲しい．

なお，突貫工事で作ったため至らない点が多々あると思われる．
バグや改善提案については GitHub 上で issue を投げてくれると，時間があるときに対応するかもしれない，
もちろんメンテしてくれる人がいれば大歓迎である．


\section*{謝　　　辞}
本研究は○○○○○○の支援のもと実施された．

\bibliographystyle{plainnat}
\bibliography{reference}

\end{document}
